\documentclass[12pt]{article}
\usepackage{amsmath,amssymb}
\begin{document}
\section{Chapter 6.1: Set Theory}
\textbf{Problem 5.13:} Prove by Mathematical induction $\forall n \in \mathbb{Z}^+ \cup \{0\}, x,y \in \mathbb{Z}, x \neq y, x - y | x^n - y^n$

\textbf{Proof} Suppose $x,y \in \mathbb{Z} \vee x \neq y.$ Let the predicate $P(n)$ be the statement $x - y | x^n - y^n$.

\textbf{Basis Step} Show that $P(0)$ is true. $P(0)$ is true because $x^0 - y^0 = 1 - 1$ and $x - y | 0.$

\textbf{Inductive Step} Show that $\forall k \in \mathbb{Z}^+ \cup \{0\}, P(k) \rightarrow P(k + 1)$

Let $k \in \mathbb{Z}^+ \cup \{0\}$ and suppose $x - y | x^k - y^k$.
We must show that \[x - y | x^{k + 1} - y^{k + 1}\] .
By rewriting the inductive hypothesis using the definition of divisiblity, we have:

\[\exists r \in \mathbb{Z}, x^k - y^k = (x- y)r\]
Then: \[x^{k + 1} - y^{k + 1}\]
= \[x^{k + 1} - xy^k + xy^k - y^{k + 1}\]
= \[x(x^k - y^k) + y^k(x - y)\]
= \[x(x - y)r + y^k(x - y)\] %by the inductive hypothesis
= \[(x - y)[xr + y^k]\]

Now \[x(r + y^k) \in \mathbb{Z}\] because products and sums of integeres are integers. Thus, \[x - y | x^{k + 1} - x^{k + 1}\] .

Let the set \[A = \{x \in \mathbb{Z} | x = 5a + 2, a \in \mathbb{Z}\}\]
\[B = \{y \in \mathbb{Z} | y = 10b - 3, b \in \mathbb{Z}\}\]
\[C = \{z \in \mathbb{Z} | z = 10c + 7, z \in \mathbb{Z}\}\]

A. Is \( A \subset B? \) . No, because  2 = 5 * 0 + 2, but \(2 \not\subset B\) because if \(2 \in B\) then \(\exists b \in \mathbb{Z}\) such that \(2 = 10b - 3\) . Then \(10b = 5, b = \frac{1}{2} \) which is not an integer. Counterexample found.

B. Is \( B \subset A? \) . Yes. Proof: Suppose \( y \in B. \) [*We must show that \(y \in A\) . By definition of A, this means that we must show \(y = 5(\text{some int}) + 2\)].
By definition of B, \(\exists b \in \mathbb{Z}, y = 10b - 3.\)
Let \(a = 2b - 1\) . Then \(a \in \mathbb{Z}\) and \(5a + 2\) =
\[5(2b - 1) + 2 = 10b - 5 + 2 = 10b - 3 = y\]

*To prove A = B, prove both are subsets of each other*
Let the universal set be \(\mathbb{R}\)
\[A \cup B = \{x \in \mathbb{R} | 6 < x \leq 8 \}\]
\[A \cap B = \{x \in \mathbb{R} | -1 < x \leq 0 \}\]
\[A^c = \{x \in \mathbb{R} | -3 > x \text{ or } x > 0 \}\]
\[A \cap C = \emptyset\]

*Draw the venn diagram to describe \(A \cap B  \neq \emptyset, B \cap C = \emptyset, A \cap C = \emptyset, A \not\subset B, C \not\subset B\)*

Questions:
Is \( 0 \in \emptyset\)? No
\(\emptyset \neq \{\emptyset\}\)? Yes.
\(\emptyset \in \{\emptyset\}\)
\(\emptyset \notin \emptyset\)

EXAM QUESTION
\[\forall i \in \mathbb{Z}^+, R_i = \{x \in \mathbb{R} | 1 \leq x \leq 1 + \frac{1}{i}\} = [1, 1 + \frac{1}{i}]\]

\[\cup_{i = 1}^{4} R_i = R_1 \cup R_2 \cup R_3 \cup R_4 = \cup [1, 1 + \frac{1}{1}] \cup [1, 1 + \frac{1}{2}] \cup [1, 1 + \frac{1}{3}] \cup [1, 1 + \frac{1}{4}] = [1,2]\]
\[\cap_{i = 1}^{4} R_i = R_1 \cap R_2 \cap R_3 \cap R_4 = \cap [1, 1 + \frac{1}{1}] \cap [1, 1 + \frac{1}{2}] \cap [1, 1 + \frac{1}{3}] \cap [1, 1 + \frac{1}{4}] = [1, 1 + \frac{1}{4}]\]

Is \(\{\{w, x, v\}, \{w,y,q\}, \{p,z\}\}\) a partition of \(\{p,q,u,v,w,x,y,z\}\)? Yes.
\end{document}
